\paragraph{Extraction.}
The support at every $p\mid m$ has $v_pQ_B=1$.  All nonzero
indices are integral and divisible by $m$.  Repeat the preceding
construction for $Q_B'=Q_B/m$ with the same finite weights and the
Gaussian for $Q_B'$.  This quadratic form is integral on the
support.  Away from $m$ it differs by a unit, so the good-prime
self-duality argument remains valid.  Its determinant squareclass
is $m^{-3}\equiv m$, giving the extra Dirichlet character
$(m/\cdot)$.  Hence
\[
 g_M(\tau)=\sum_n H_M(mn)q^n\in
      M_{3/2}(\Gamma_0(K),\psi).
\]
Here $K$ has only the stated prime support, and
\[
 \psi=\eta_B(m/\cdot),\qquad
 \eta=\eta_B\chi_{h_*},\qquad
 \psi/\eta=\chi_{-1}. \tag{O.5}
\]
This directly constructs the extracted form without a
half-integral-weight extraction operator.
The inverse in this classical passage can also be checked before
any character notation is introduced. Put
$g_\tau=n(x)m(\sqrt y)$ for $\tau=x+iy$, so $g_\tau i=\tau$.
For $\gamma\in\operatorname{SL}_2(\mathbb Q)$ write
$\gamma g_\tau=g_{\gamma\tau}k_\infty$ with
$k_\infty\in\operatorname{SO}_2(\mathbb R)$. In the adelic group,
\[
 (g_{\gamma\tau},1_f)
   =\gamma\,(g_\tau,\gamma_f^{-1})(k_\infty^{-1},1_f).
\]
This equality also holds with consistent metaplectic lifts.
The global rational lifts are compatible for the two forms:
\cite[II.30, Proposition~5]{weil} gives
$\prod_v\gamma_v(Q)=\prod_v\gamma_v(aQ)=1$ for rational $Q,a$,
and hence the local normalizing factors in $W_Q$ and $M_Q(a)$
have product one. The product formula also gives
$\prod_v|\det(2A)|_v^{1/2}=1$, so the product of the stated
self-dual measures is the measure for which the rational lattice
has covolume one. Apply the rational invariance of
\cite[III.41, Theorem~6]{weil} to the displayed adelic identity.
The archimedean Gaussian for either positive ternary form has
the same weight $3/2$; therefore its factor from $k_\infty^{-1}$
is the same as that for $\theta^3$. The relative finite factor
is evaluated at $\gamma_f^{-1}$, whose upper-left entry is $d$.
Thus the relative classical character is
$\prod_p\alpha_p(d)(d,\det Q)_p$, over the prescribed primes.
This proves the inverse-idele convention directly. It compares
the transformation identities and does not divide by a value
of a theta function.

For the extracted form there is likewise no measure or degree
normalization in its coefficients. With exactly the same finite
weights, replace only the quadratic form and archimedean Gaussian
by $Q_B/m$ and $e^{-2\pi Q_B/m}$. The summand indexed by $x$ now
has exponent $Q_B(x)/m$. The original support makes $Q_B(x)$ an
integer with valuation one at every prime dividing $m$, so these
exponents are nonnegative integers. Coefficient $n$ is exactly
the original sum over $Q_B(x)=mn$, namely $H_M(mn)$.
The self-dual measure changes by $|m|_v^{-3/2}$ locally, but a
theta coefficient is a discrete sum and contains no Haar
measure; this change is used only in the Fourier operators.
Outside the original support it is a unit rescaling of the same
unimodular local form, and at primes dividing $m$ the
lower-unipotent depth can be increased within the existing
support. The determinant changes by $m^{-3}$, whose squareclass
is $m$, so (O.5c) supplies exactly $(m/\cdot)$.
In particular no new level prime, factor of $m$, or dyadic
root of unity is introduced by extraction.


